\section[The duality of W\_-1(sl\_4, f\_sub) and L\^{}N=2\_c=-15]{The duality of $\boldsymbol{\mathcal{W}_{-1}(\mathfrak{sl}_4, f_{\rm sub})}$ and $\boldsymbol{L^{N=2}_{c=-15}}$} \label{duality}


As we have seen in Section \ref{KS_gen}, Kazama--Suzuki duality between $L^{N=2}_c$ and $\mathcal W_k(\mathfrak{sl}_4, f_{\rm sub})$ can only occur if $k=-1$ and $c=-15$ or $k=-\frac{7}{3}$ and $c=1$.
In this section, we consider the case~${k=-1}$. Then the central charge of $\mathcal{W}^{k}(\mathfrak{sl}_4, f_{\rm sub})$ is $c = -15$. According to Lemma \ref{sing_G^+}, \smash{$(G^+)^2$} is singular vector in $\mathcal{W}^{k}(\mathfrak{sl}_4, f_{\rm sub})$, hence \smash{$(G^+)^2 = 0$} in $\mathcal{W}_{k}(\mathfrak{sl}_4, f_{\rm sub})$.

 We will show that $\mathcal{W}_{-1}(\mathfrak{sl}_4, f_{\rm sub})$ is the Kazama--Suzuki dual of the $N=2$ superconformal vertex algebra of central charge $-15$.



\subsection[Embedding of L\^{}N=2\_c=-15 into W\_-1(sl\_4, f\_ ub)]{Embedding of $\boldsymbol{ L^{N=2}_{c=-15}}$ into $\boldsymbol{\mathcal{W}_{-1}(\mathfrak{sl}_4, f_{\rm sub}) \otimes \F_{-1}}$}

First, we will show that the $N=2$ superconformal algebra $L^{N=2}_{c}$ for $c=-15$ can be realized as a subalgebra of $\mathcal{W}_{-1}(\mathfrak{sl}_4, f_{\rm sub}) \otimes \F_{-1}$, where $\F_{\pm 1}$ is the vertex algebra associated to the rank one lattice $L= {\mathbb{Z}}\varphi^{\pm}$, with $\big\langle \varphi^{\pm}, \varphi^{\pm} \big\rangle = \pm 1$ (cf.\ Section \ref{F_{-1}}).

 In this section, we shall use that fact that the universal superconformal $N=2$ vertex algebra~$V^{N=2}_{c}$ is simple for $c=-15$, since it is the Kazama--Suzuki dual of the universal affine vertex algebra $V^{k=-5/3}(\mathfrak{sl}_2)$ (which is also simple by \cite{GK}).

 Note that the maximal ideal in $\mathcal W^k(\mathfrak{sl}_4, f_{\rm sub})$ is invariant under the automorphism which maps $G^+$ to $G^-$. Since \smash{$(G^+)^2$} is a singular vector in $\mathcal{W}^{-1}(\mathfrak{sl}_4, f_{\rm sub}) $, we conclude that $(G^-)^2$ also belongs to the maximal ideal of $\mathcal{W}^{-1}(\mathfrak{sl}_4, f_{\rm sub}) $. Now we want to show that the maximal ideal is generated by these two vectors.
Let \smash{$I = \bigl\langle (G^+)^2, (G^-)^2 \bigr\rangle$} and define
\[
 \widetilde W = \mathcal{W}^{-1}(\mathfrak{sl}_4, f_{\rm sub})/I.
 \]

Let $H^{\perp} = J + \varphi^-$. Then for $n \ge 0$, we have
\[
H^{\perp}(n) H^{\perp} = \frac{1}{4} \delta_{n,1}\mathbbm{1}.
\]
Let $M_{H^{\perp}} (1)$ be the Heisenberg vertex algebra generated by $H^{\perp}$, and $M_{H^{\perp}} (1, s)$ the irreducible highest weight $M_{H^{\perp}} (1)$-module on which
$H^{\perp}(0) \equiv s \operatorname{Id}$.

\begin{Proposition} \label{ulaganje-1}\qquad
\begin{itemize}\itemsep=0pt
\item[$(1)$]There is a vertex algebra homomorphism $\Phi\colon L^{N=2}_{c=-15} \rightarrow \widetilde W \otimes \F_{-1}$ given by
	\begin{gather*}
		E = \frac{2}{3} G^+ \otimes e^{\varphi^-}, \qquad
		F = - G^- \otimes e^{-\varphi^-},\qquad
		H = - 4 J - 5 \varphi^-, \\
		T = L- 2{:}J J {:}-4{:}J\varphi^-{:} -\frac{5}{2}{:}(\varphi^-)^2{:}.
	\end{gather*}
	\item[$(2)$]
	Let $H^{\perp} = J + \varphi^-$. Then
\[
 \widetilde W \otimes \F_{-1} \cong \bigoplus_{n \in {\mathbb{Z}}} \sigma^n \bigl(L^{N=2}_{c=-15}\bigr) \otimes M_{H^\perp} (1, n).
 \]
	\item[$(3)$] We have
\[
\operatorname{Im} (\Phi) \cong \operatorname{Com} (M_{H^{\perp}} (1), \mathcal{W}_{-1}(\mathfrak{sl}_4, f_{\rm sub}) \otimes \F_{-1}) \cong L^{N=2}_{c=-15} .
\]
\item[$(4)$] $\widetilde W = \W_{-1} (\mathfrak{sl}_4, f_{\rm sub})$, i.e., $I$ is the maximal ideal in $\W^{-1}(\mathfrak{sl}_4, f_{\rm sub})$.
\end{itemize}
	
\end{Proposition}

\begin{proof}


It is easy to check that
\begin{gather*}
 H(0) E = E, \qquad H(0) F =-F, \qquad H(n) E = H(n) F =0, \qquad n >0, \\
 H(1) H = - 5 {\mathbbm 1} = \frac{c}{3} {\mathbbm 1},\qquad H(n) H =0, \qquad n >1.
\end{gather*}
 Let $L^{\perp} = L- \frac{2}{5} {:}J J{:}$ (cf.\ Appendix \ref{appendix-b}). We get that
\[
 T = L^{\perp} -\frac{1}{10} {:}H H{:}.
 \]
This implies that $T$ is a Virasoro vector of central charge $c=-15$. Clearly, we have that $H$ is a primary vector for $T$ of conformal weight $1$.


 Next we notice that $\bigl(G^\pm \bigr)^2 =0$ in $\widetilde W$. Direct calculation shows that
 \begin{gather*}
 E_{(0)} E = \frac{4}{9} {\bigl( G^+ \bigr)}^2 \otimes e^{2 \varphi^-} =0, \qquad F_{(0)}F = {( G^-)}^2 \otimes e^{-2 \varphi^-} =0, \\
 E_{(n)} E = F_{(n)} F =0, \qquad n\in {\mathbb Z}_{\ge 0},
\end{gather*}
and
\begin{gather*}
	E_{(2)} F = -\frac{2 (2 + k)(5 + 2k)(8 + 3k)}{3} {\mathbbm 1} =(-10) {\mathbbm 1} = \frac{2 c}{3} {\mathbbm 1}, \\
E_{(1)} F = -\frac{ 8(2 + k)(5 + 2k)}{3} J - \frac{2 (2 + k)(5 + 2k)(8 + 3k)}{3} \varphi^- = -8 J - 10 \varphi^- = 2 H, \\
 E_{(0)} F = - \frac{2 (2 + k)(5 + 2k)(8 + 3k)}{3}\cdot\frac{1}{2}\bigl({:}( \varphi^-)^2{:}+\partial\varphi^-\bigr)-\frac{ 8(2 + k)(5 + 2k)}{3} {:}J \varphi^-{:} \\
 \phantom{E_{(0)} F =}{}+\frac{2(k+2) \bigl((k+4)L-6{:}J^2{:}-2(2k+5) \partial J\bigr)}{3} = 2T + \partial H.
\end{gather*}

The above relations show that the even fields $H$, $ T$ and the odd fields $E$, $F$ satisfy the $\lambda$-bracket for the $N=2$ superconformal vertex algebra of central charge $c=-15$. Since \smash{$V^{N=2} _{c=-15}$} is simple, we conclude that $H$, $T$, $E$, $F$ generate a vertex subalgebra of $\widetilde W \otimes \F_{-1}$ isomorphic to~$L^{N=2}_{c=-15}$. This proves the assertion~(1).

 Let \smash{$ \overline W = \operatorname{Ker} _{ \widetilde W\otimes F_{-1}} H^{\perp}(0)$}. Then $\overline{W}$ is a vertex algebra which contains $ \operatorname{Im} (\Phi)\otimes M_{H^{\perp}} (1)$.



Let us prove the following claim.
\begin{Claim}\label{Claim1}
 $ \overline W $ is a simple vertex algebra generated by $\big\{ E,F,T,H, H^{\perp} \big\}$.
\end{Claim}

Let $U$ be the vertex subalgebra of $\overline W $ generated by $ \big\{E,F,T,H, H^{\perp} \big\}$. Then $U \cong \operatorname{Im} (\Phi) \otimes M_{H^{\perp}} (1)$. Since \smash{$\operatorname{Im} (\Phi) \cong L^{N=2}_{c=-15}$}, we conclude that
\smash{$U\cong L^{N=2}_{c=-15} \otimes M_{H^{\perp}} (1)$}, and it is therefore simple.

 For each $n \in {\Z}$, we consider the $U$-module $U^{(n)} = U. e^{n \varphi^-}$. One can show that
 $ U^{(n)} $ is isomorphic to the simple $U$-module obtained by the simple current construction
\[
 \bigl(U^{(n)}, Y^{(n)} (\cdot, z)\bigr): = (U , Y(\Delta(n \varphi^-,) \cdot, z)).
 \]


 Note that $ \varphi^- = -H - 4 H^{\perp}$, which implies that
 \[% \label{delta-alpha}
 \Delta(\varphi^-, z) = \Delta( -H, z) \Delta\bigl(- 4 H^{\perp}, z\bigr).
 \]
 \begin{itemize}\itemsep=0pt
\item For a $ \operatorname{Im} (\Phi)$-module $M$, by applying the operator $\Delta(-nH, z)$, we get the module $\sigma^{n}( M)$.
\item Applying the operator $ \Delta\bigl( -4n H^{\perp}, z\bigr)$ on
 $ M_{H^{\perp}}(1)$, we get the $ M_{H^{\perp}}(1)$-module
$ M_{H^{\perp}} (1, n )$.
\end{itemize}

We get
$
U^{(n)}= \sigma^{n}( \operatorname{Im} (\Phi)) \otimes M_{H^{\perp}} (1,n )$.
Note that $ H^{\perp} (0) \equiv n \operatorname{Id} $ on $U^{(n)}$.
 Using the construction of H. Li from \cite{Li-ext}, we get that
\[%\label{dec-u1}
 \mathcal U = \bigoplus_{n \in \mathbb{Z}} U^{(n)}
 \]
is a vertex algebra and hence a vertex subalgebra of $ \widetilde W \otimes \F_{-1}$. But it is not hard to see that $ \mathcal U$ contains all generators of $ \widetilde W \otimes \F_{-1}$.

 Indeed, $e^{\varphi^-}\in U^{(1)}$, $e^{-\varphi^-}\in U^{(-1)}$ and
 \[
 G^+ = \frac{3}{2} E_{(-2)}e^{-\varphi^-}, \qquad G^- =- F_{(-2)}e^{\varphi^-}.
 \]


Since $G^+$ and $G^-$ (weakly) generate $ \widetilde W$, it follows that
$
 \mathcal U = \widetilde W \otimes \F_{-1}$.






This proves that $ U= \overline W \cong \operatorname{Im} (\Phi) \otimes M_{H^{\perp}} (1) $. Therefore, $\overline W$ is a simple vertex algebra. This proves Claim~\ref{Claim1}.

Now we get
\[
 \widetilde W \otimes \F_{-1} = \bigoplus_{n \in \mathbb{Z}} U^{(n)} = \bigoplus_{n \in \mathbb{Z}} \sigma^n \bigl(L^{N=2}_{c=-15}\bigr) \otimes M_{H^\perp} (1, n).
 \]
This implies the assertions (2) and (3).

As the right-hand side of decomposition in (2) is simple, it follows that $\widetilde W $ must be simple and therefore isomorphic to the unique simple quotient $\mathcal{W}_{-1}(\mathfrak{sl}_4, f_{\rm sub})$. Hence $I$ is the maximal ideal in $\mathcal{W}^{-1}(\mathfrak{sl}_4, f_{\rm sub})$.
This proves assertion (4).
\end{proof}

 Let $M_H(1)$ be the Heisenberg subalgebra of $L^{N=2}_{c=-15}$ generated by the Heisenberg field $H(z)$ and $M_J(1)$ the Heisenberg subalgebra of $\mathcal{W}_{k}(\mathfrak{sl}_4, f_{\rm sub}) $ generated by the Heisenberg field $J(z)$.
From the proof of Proposition~\ref{ulaganje-1}, we have the following corollary.

	\begin{Corollary} \label{pf} $\operatorname{Com} \bigl(M_H(1) , L^{N=2}_{c=-15} \bigr) \cong \operatorname{Com} (M_J(1), \mathcal{W}_{-1}(\mathfrak{sl}_4, f_{\rm sub}) )$. \end{Corollary}
	
\subsection[Embedding of W\_k(sl\_4, f\_sub) into L\^{}N=2\_c=-15]{Embedding of $\boldsymbol{\mathcal{W}_{k}(\mathfrak{sl}_4, f_{\rm sub})}$ into $\boldsymbol{L^{N=2}_{c=-15} \otimes \F_{1}}$ }

It remains to show that $\mathcal{W}_{-1}(\mathfrak{sl}_4, f_{\rm sub}) $ can be realized as a subalgebra of the tensor product of the $N=2$ superconformal algebra $L^{N=2}_{c}$ for $c=-15$ and the lattice vertex algebra $\F_1$.

Let $\overline{H} = H - \varphi^+$. Then for $n \ge 0$ we have $\overline H(n) \overline H = -4 \delta_{n,1}\mathbbm{1}$.
Let $M_{\overline{H}} (1)$ be the Heisenberg vertex algebra generated by $\overline H$, and $M_{\overline{H}} (1, s)$ the irreducible highest weight $M_{\overline{H}} (1)$-module on which
$\overline H(0) \equiv s \operatorname{Id}$.





\begin{Proposition} \label{ulaganje-2}
\quad
\begin{itemize}\itemsep=0pt
	\item[$(1)$] There is a vertex algebra homomorphism $\Phi^{\rm inv} \colon \mathcal{W}^{-1}(\mathfrak{sl}_4, f_{\rm sub}) \rightarrow L^{N=2}_{c=-15} \otimes \F_{1}$ given by
	\begin{gather*}
		G^+ = E\otimes e^{\varphi^+}, \qquad
		G^- = \frac{3}{2} F \otimes e^{-\varphi^+}, \qquad
		J = - \frac{1}{4} H + \frac{5}{4} \varphi^+, \\
		L = T + \frac{1} {10} {:}H H{:}+ \frac{2}{5} {:}\left( - \frac{1}{4} H + \frac{5}{4} \varphi^+\right)^2{:}, \qquad
		W = -\frac{3}{2} W^{N=2}_{c=-15}.
	\end{gather*}
	
\item[$(2)$] $\operatorname{Im} \bigl(\Phi^{\rm inv}\bigr) $ is isomorphic to the simple vertex algebra $ \mathcal{W}_{-1}(\mathfrak{sl}_4, f_{\rm sub})$.

\item[$(3)$]
 As a $ \mathcal{W}_{-1}(\mathfrak{sl}_4, f_{\rm sub})\otimes M_{\overline{H}} (1)$-module
\[
L^{N=2}_{c=-15} \otimes \F_{1} \cong \bigoplus_{n\in {\Z} }
 \psi^{-n} ( \mathcal{W}_{-1}(\mathfrak{sl}_4, f_{\rm sub})) \otimes M_{\overline{H}} (1, n).
 \]
 \item[$(4)$] We have
\[
 \operatorname{Com} \bigl(M_{\overline{H}}(1), L^{N=2}_{c=-15} \otimes \F_{1}\bigr) \cong \mathcal{W}_{-1}(\mathfrak{sl}_4, f_{\rm sub}).
 \]
 \end{itemize}
	
\end{Proposition}
\begin{proof}
(1)
 First we notice that
$L = T^{\perp} + \frac{2}{5}{:}JJ{:}$,
which implies that $L$ is a Virasoro vector of central charge $c=-15$. Note that
$
 L^{\perp} = L- \frac{2}{5}{:}JJ{:} = T^{\perp}$.
Clearly, $J$ is a primary vector of conformal weight $1$ with respect to $L$.
Moreover, since
$L = T + \frac{2}{5}{:}JJ{:} + \frac{1}{10} {:}H H{:} $, we conclude that $G^{\pm}$ are primary vectors of conformal weight $2$.
Direct calculation shows that
\begin{gather*}
 G^+(3) G^- = \frac{3}{2}(-E_{(2)}F) = \frac{3}{2}\left(- \frac{2 c}{3}\right) \mathbbm{1} =15\mathbbm{1} =(2 + k)(5 + 2k)(8 + 3k)\mathbbm{1}, \\
 G^+(2) G^- =\frac{3}{2}\left(-E_{(1)}F- \varphi^+(-1) E_{(2)}F\right) \\
 \phantom{G^+(2) G^- }{} = \frac{3}{2}\left(-2H -\frac{2c}{3} \varphi^+\right) = 12\left( - \frac{1}{4} H + \frac{5}{4} \varphi^+ \right) =12J = 4(2 + k)(5 + 2k)J.
 \end{gather*}


 We should show that
\[ G^+(1) G^-
 = - 3 L + 6 {:} JJ {:} + 6 \partial J = -3 L^{\perp} + \frac{24}{5} J^2 + 6 \partial J.
 \]
 Indeed
\begin{align*}
G^+(1) G^- &= \frac{3}{2}\left(-E_{(0)}F-\varphi^+(-1) E_{(1)}F-\frac{1}{2}\bigl({:}\bigl(\varphi^+\bigr)^2{:}+\partial\varphi^+\bigr)E_{(2)}F \right) \\
&= \frac{3}{2}\left( -2 T -\partial H -2{:}H\varphi^+{:} - \frac{1}{2}\bigl({:}\bigl(\varphi^+\bigr)^2{:}+\partial\varphi^+ \bigr)\frac{2c}{3}\right)\mathbbm{1} \\
& = -3 L^{\perp}+ \frac{3}{10} {:}H^2{:}- \frac{3}{2} \partial H -3 {:}H \varphi^+ {:} + \frac{15}{2} \bigl({:}\bigl(\varphi^+\bigr)^2{:} +\partial \varphi^+\bigr) \\
& = -3 L^{\perp} + \frac{24}{5} {:}J ^2{:} + 6\partial J.
 \end{align*}

 Next, we need to show that
 \begin{align}
G^{+}(0) G^{-} ={}&(k+2)\left(W + \frac{8(11k+32)}{3(3k+8)^2}{:}J^3{:}-\frac{4(k+4)}{3k+8}{:}L J{:} + 6{:}\partial JJ{:}\right. - \frac{k+4}{2}\partial L \nonumber\\
 &+ \left. \frac{4(3k^2+17k+26)}{3(3k+8)}\partial^2J\right) \nonumber \\
={}& W+ \frac{32}{25}{:}J^3{:}-\frac{12}{5}{:}JL^{\perp}{:}+\frac{24}{5} {:}\partial JJ{:}-\frac{3}{2}\partial L^{\perp}+2\partial^2J.\label{GpGm-OPE}
\end{align}

We have that
 \begin{align}
G^+(0) G^-={}& \frac{3}{2}\left(-E_{(-1)}F -\varphi^+(-1) E_{(0)}F - E_{(1)}F\cdot \frac{1}{2}\bigl({:}\bigl(\varphi^+\bigr)^2{:}+\partial \varphi^+\bigr) \right. \nonumber\\
&\left. - E_{(2)}F\cdot \frac{1}{6}\bigl({:} \bigl(\varphi^+ \bigr)^3 {:}+3{:}\varphi ^+ \partial \varphi^+{:} + \partial ^2 \varphi^+\bigr) \right) \nonumber\\
={}& \frac{3}{2}\left(-E_{(-1)}F -\varphi^+(-1) \left(2L^{\perp} - \frac{1}{5}{:}H^2{:} +\partial H\right) - 2H\cdot\frac{1}{2}\bigl({:}\bigl(\varphi^+\bigr)^2{:}+\partial \varphi^+\bigr) \right. \nonumber\\
&\left.- \frac{2c}{3}\cdot \frac{1}{6}\bigl({:}\bigl(\varphi^+\bigr)^3{:} + 3{:}\varphi^+ \partial \varphi^+{:} + \partial ^2 \varphi^+\bigr) \right) .  \label{GpGm-rel}
\end{align}




As the parafermionic subalgebras of $\mathcal W_{-1} (\mathfrak{sl}_4, f_{\rm sub})$ and $L^{N=2}_{c=-15} $ coincide (cf.\ Corollary \ref{pf}), it follows that the field $W$ coincides (up to normalization) with the parafermionic generator~$W^{N=2}_{c=-15}$, that is,
\begin{align}
	W	&= \nu \left({:}E F{:} - \partial L^{\perp} - \frac{3}{2c}{:}H\partial H{:} -\frac{6}{c}{:}HL^{\perp}{:}- \frac{1}{3}\partial^2H + \frac{3}{c^2}{:}H^3{:} \right) \nonumber \\
	&= \nu \left({:}E F{:} - \partial L^{\perp}+ \frac{1}{5}{:}\partial H H {:} + \frac{2}{5}{:}HL^{\perp}{:} -\frac{1}{3}\partial^2H - \frac{1}{75}{:}H^3{:} \right). \label{W}
\end{align}
Setting $\nu = -\frac{3}{2}$ and substituting \eqref{W} into \eqref{GpGm-OPE}, we obtain that \eqref{GpGm-rel} = \eqref{GpGm-OPE}.
This proves the assertion (1).

(2) Let us prove that $\operatorname{Im} \bigl(\Phi^{\rm inv}\bigr) $ is simple. Let \smash{$\overline{W} = \operatorname{Ker}_{ L^{N=2}_{c=-15} \otimes \F_{1}} \overline{H} (0)$}. It is clear that $\overline{W}$ is a simple vertex algebra which contains $ \operatorname{Im} \bigl(\Phi^{\rm inv}\bigr) \otimes M_{\overline{H}} (1)$. The simplicity of $ \operatorname{Im} \bigl(\Phi^{\rm inv}\bigr) $ follows from the following claim.

\begin{Claim}\label{Claim2}
$\overline W$ is generated by $\big\{ G^+, G^-, J, L, W, \overline{H} \big\}$.
\end{Claim}

Proof of Claim~\ref{Claim2} is analogous to the proof of Claim~\ref{Claim1} in Proposition \ref{ulaganje-1}. Let $U$ be the vertex subalgebra of $\overline W$ generated by $\big\{ G^+, G^-, J, L^U, W, \overline{H} \big\}$. Then clearly $U \cong \operatorname{Im} \bigl(\Phi^{\rm inv}\bigr) \otimes M_{\overline{H}} (1)$.

Let $ U^{(n)} $ be the $U$-module obtained by the simple current construction
\[ \bigl(U^{(n)}, Y^{(n)} (\cdot, z)\bigr): = \bigl(U , Y\bigl(\Delta\bigl(n \varphi^+\bigr) \cdot, z\bigr)\bigr).
\]

As in Proposition \ref{ulaganje-1}, from the formula
\[% \label{delta-alpha}
\Delta\bigl(n\varphi^+, z\bigr) = \Delta( nJ, z) \Delta\left(\frac{n}{4} \overline{H}, z\right)
\]
we get
$U^{(n)}= \psi^{-n}\bigl( \operatorname{Im} \bigl(\Phi^{\rm inv}\bigr)\bigr) \otimes M_{\overline{H}} (1,n )$.


The rest of the proof follows analogously.
\end{proof}

We have proved the following.

\begin{Theorem} The vertex algebra $\mathcal{W}_{-1}(\mathfrak{sl}(4), f_{\rm sub})$ is the Kazama--Suzuki dual of the $N=2$ superconformal vertex algebra $L^{N=2}_{c=-15}$.
\end{Theorem}
